Du haut d'un pont, un enfant laisse tomber une pierre sans vitesse initiale. La pierre touche l'eau \qty{\v{h}}{m} plus bas. On néglige la résistance de l'air et on prend $g = \qty{9.81}{m/s^2}$.

\begin{enumerate}[label=\alph*)]
    \item Combien de temps la pierre met-elle pour atteindre l'eau ?
    \item Quelle est la vitesse de la pierre juste avant de toucher l'eau ?
    \item À quelle hauteur au-dessus de l'eau la pierre se trouve-t-elle \qty{\v{t1}}{s} après avoir été lâchée ?
    \item Le son du plouf remonte jusqu'à l'enfant à \qty{340}{m/s}. Combien de temps après avoir lâché la pierre l'enfant entend-il le plouf ?
\end{enumerate}
